\section{Nine dimensions: winding labels and the exceptional groups}
\label{sec:ninedim}
On a circle the two heterotic strings are one theory, so the label
descriptions of Sections~\ref{sec:gaugefieldscalars}
and~\ref{sec:threelabel} can be compared. This section records the
comparison. Details and the exact-arithmetic script are in
blackboard~015. The lattice statements are computed. The diagram of the
lattice and the list of points of maximal symmetry are
known~\cite{FFGNP}. The reading in terms of flavor, color and charge is
conditional, and remarks on the type~I$'$ description are quoted from
memory.

\paragraph{Setting.}
Charges form the even self-dual lattice of signature $(17,1)$. In
$SO(32)$ variables a vector is $v=(\pi;n,w)$, with $\pi$ the sixteen
label charges, $n$ the momentum and $w$ the winding number, and
$v^2=\pi^2+2nw$. A vacuum is a Wilson line $A$ and a radius $R$,
combined in $t=(-A;-(A^2/2+r),1)$ with $r=R^2/\alpha'$. A vector with
$v^2=2$ is a state with
\begin{equation}
\alpha' M^2=\frac{(v\cdot t)^2}{r}
\label{eq:ninedmass}
\end{equation}
in units $\hbar=c=1$ and with physical label charges $\pi+wA$; the
massless ones are the charged gauge bosons. A set of nineteen simple
roots is $g_1=(0;-1,-1)$, $a_0=(-e_1-e_2;1,0)$,
$a_i=(e_i-e_{i+1};0,0)$ for $i=1,\dots,15$,
$a_{16}=(e_{15}+e_{16};0,0)$ and $g_2=((-\tfrac12)^{16};1,-1)$. Two of
them carry winding. Deleting two nodes of their diagram gives 44 points
of maximal symmetry up to its left--right symmetry, and a direct
enumeration of roots agrees with the diagram at each.

\paragraph{Exceptional groups.}
With $m$ labels at Wilson-line value $1/2$, the others at $0$ and
$r=(8-m)/8$, the gauge group is $E_{m+1}\times SO(32-2m)$, where
$E_4=SU(5)$ and $E_5=SO(10)$. The exceptional factor consists of the
$SO(2m)$ of the $m$ labels at zero winding, its spinor at winding
$\pm1$ and, for $m=6$ and $7$, a singlet or the vector at winding
$\pm2$. For $m=7$,
\begin{equation}
248=14+64+92+64+14\qquad(w=-2,\dots,2),
\end{equation}
and the winding-one states are the subsets of the seven labels of even
size. If the seven labels are $(u,c)$ times color and the sixteenth
label, three colored labels come with the sixteenth. The color singlets
at winding one are then the empty set, the four flavor-symmetric states
$uuu$, $uuc$, $ucc$, $ccc$ with the sixteenth label, and the state with
all six colored labels.

\paragraph{The nine labels of $E_8$.}
An $SU(9)$ inside this $E_8$ has seven $SO(32)$ labels and two labels
that differ from them by winding: the nine labels carry winding
$(\tfrac13)^7,-\tfrac23,-\tfrac53$. The three-label multiplet splits as
$84=35+21+21+7$ at winding $1,0,-1,-2$. The 21 at zero winding are
$SO(32)$ two-label states, among them the fifteen diquark-like states of
the six colored labels. A diquark-like state of $SO(32)$ is thus a
three-label state of $E_8$ whose third label is a winding label. Other
duality frames grade $E_8$ differently: among 6000 frames generated by
reflections its zero-winding part is $SO(14)$, $SU(8)$, $SO(12)$,
$SU(7)\times SU(2)$, $SU(7)$ or $SO(10)\times SU(3)$. With $SU(8)$,
eight of the nine labels are $SO(32)$ labels of that frame. This
replaces by a computation the argument at the end of
Section~\ref{sec:threelabel}.

\paragraph{Color.}
In any frame of $SO(32)$ type the roots of $E_8$ at zero winding embed
in the roots of $D_{16}$, which has no subsystem of type $E$. For nine
labels read as $3_f\otimes3_c$, every color-blind direction is orthogonal
to the roots of an $E_6$. Hence the winding number never commutes with
the color group of Section~\ref{sec:threelabel}, and the nine labels
are never all $SO(32)$ labels: the decuplet reading of $E_8$ has no
expression in $SO(32)$ labels. The color group of
Section~\ref{sec:gaugefieldscalars} behaves differently. At the
$E_8\times SO(18)$ point it sits at zero winding and acts on the nine
labels as $3+3+1+1+1$. Its singlets in the 248 number 22 and form
$G_2\times SU(3)$. The $SU(3)$ acts on the three colorless labels, and
the long roots of $G_2$ are $uuu$, $ccc$, the state of the three
colorless labels and their conjugates.

\paragraph{Vacua that respect color and charge.}
Of the 44 points, 24 keep the color group of
Section~\ref{sec:gaugefieldscalars} at zero winding. For 11 of them the
Wilson line is constant on the six $(u,c)$ labels and on the nine
$(d,s,b)$ labels (Table~\ref{tab:ninedim}). In that class exceptional
symmetry arises only as $E_7$ or $E_8$ and only on the up-type block,
with the sixteenth label deciding between the two.

\begin{table}[ht]
\centering
\caption{Points of maximal gauge symmetry of the heterotic string on a
circle whose Wilson line depends only on electric charge. The last
column gives the Wilson-line value of each block of labels, folded to
$[0,\tfrac12]$; $n$ is the sixteenth label.}
\label{tab:ninedim}
\begin{tabular}{@{}llll@{}}
\toprule
Blocks of labels & Gauge group & $R^2/\alpha'$ & Wilson line \\
\midrule
all sixteen & $SO(32)\times SU(2)$ & $1$ & $0$ \\
 & $SO(34)$ & $1/2$ & $0$, sign on the spinor class \\
 & $SU(16)\times SU(2)\times SU(2)$ & $1/2$ & $1/4$ \\
 & $SU(17)\times SU(2)$ & $17/49$ & $3/14$ \\
 & $SU(18)$ & $1/4$ & $1/4$ \\
$(u,c,n\,|\,d,s,b)$ & $E_8\times SO(18)$ & $1/8$ & $1/2\,|\,0$ \\
 & $E_8\times SU(9)\times SU(2)$ & $1/9$ & $1/2\,|\,1/18$ \\
 & $E_8\times SU(10)$ & $5/49$ & $1/2\,|\,1/14$ \\
$(u,c\,|\,d,s,b,n)$ & $E_7\times SO(20)$ & $1/4$ & $1/2\,|\,0$ \\
 & $E_7\times SU(10)\times SU(2)$ & $1/5$ & $1/2\,|\,1/10$ \\
 & $E_7\times SU(11)$ & $11/64$ & $1/2\,|\,1/8$ \\
\bottomrule
\end{tabular}
\end{table}

\paragraph{Massive three-label states.}
At the $E_8\times SO(18)$ point the lightest massive states have
$\alpha'M^2=1/2$ and form the spinor of $SO(18)$ at winding $\pm1$.
They are vectors of the spinor class of Section~\ref{sec:threelabel}
with one unit of winding, and they contain $\Lambda^3(9)$ of the nine
$(d,s,b)$ labels, with the ten decuplet labels as color singlets. At the
$SU(18)$ point the eighteen labels are the sixteen $SO(32)$ labels and
two winding labels. Every state there has label number divisible by
three, so the group is $SU(18)/\mathbb Z_3$, and the lightest massive
states, at $\alpha'M^2=1$, are the 816 states of $\Lambda^3(18)$ and
their conjugates. Of these, 120 have zero winding and are the
diquark-like states of $SO(32)$.

\paragraph{Strings that are points in the noncompact directions.}
Two mechanisms of this setting leave a string without extension in the
noncompact directions; blackboard~016 records them. First, a Wilson
line proportional to electric charge, $A=\tfrac32Q$, has holonomy
$(-1)^{3Q}$. A two-label state is lifted exactly when $3Q$ is odd, with
$\alpha'M^2=1/(4r)$. Of the 24 color-singlet meson-like labels the 12
neutral ones stay massless and the 12 of charge $\pm1$ are lifted. In
the type~I$'$ description a lifted state is an open string stretched
across the compact interval. For charge $-1$ of the sixteenth label this
vacuum is the $E_7\times SO(20)$ point at $r=1/4$. Second, the winding
states are D-particles in that description, and an open string between
a D-particle and a D8-brane has mixed boundary conditions in eight
directions: it has no position of its own and its only massless state
is fermionic~\cite{BGS,BGL}. At the $SU(18)$ point the 32 massless
states of winding one are single labels attached to a winding label:
fifteen quark-like color triplets and one state with the sixteenth
label and charge $-1$, with their conjugates, for the traceless charge
$\mathrm{diag}(0;q_1,\dots,q_{16};0)$. Such a state keeps the
zero-point motion of its vibration modes, a mean square displacement
$\alpha'\sum_{n\le N}1/n$ per noncompact direction, set by the slope of
the string alone and independent of its internal length. For
the fundamental string of this setting that spread is far below any
bound, so these states are point-like. The hadronic length of the
bosonic partners must then enter by another route: this is the question
of scales of the earlier sections seen from the other side.

\paragraph{What is missing.}
These vacua have sixteen supercharges and only adjoint matter. Nothing
here selects a vacuum, produces chirality or fixes the charge of
winding states, which needs one number beyond the label charges.
